% analysis/design_principles-DE.tex
\section{Entwurfsprinzipien des Alignment-Based Information Modelling}
\subsection*{Motivation}
Der AIM-Rahmen integriert Geometrie, Koordinatensysteme, Realisierung und
Optimierung in ein einheitliches Modell. Aus dieser Integration folgen
grundlegende Entwurfsprinzipien.

\subsection{Prinzip 1: Krümmungszentrierte Modellierung}
\begin{proposition}
Eisenbahn-Alignments sollen primär durch Krümmungsfunktionen \(\kappa(s)\) und
nicht durch punktbasierte Geometrie modelliert werden.
\end{proposition}
Krümmung bietet intrinsische Darstellung, unmittelbare Verbindung zur Dynamik
und natürliche Integration mit Optimierung.

\subsection{Prinzip 2: Trennung von Darstellung und Realisierung}
\begin{proposition}
Das theoretische \textbf{Alignment}-Modell muss von seiner physischen
Realisierung getrennt werden.
\end{proposition}
Der Realisierungsoperator führt Glättung, Bedingungen und Abweichungen vom
idealen Modell ein. Die Missachtung dieser Trennung erzeugt systematische
Modellierungsfehler.

\subsection{Prinzip 3: Hybride Modellierung}
\begin{proposition}
Alignment-Modelle sollen parametrische Struktur mit lokalen Korrekturtermen
verbinden.
\end{proposition}
Dies ermöglicht interpretierbare Entwurfsabsicht, robuste Behandlung
unvollkommener Daten und flexible Verfeinerung.

\subsection{Prinzip 4: Operatorbasierte Architektur}
\begin{proposition}
Alignment-Modellierung soll als Operatorkomposition formuliert werden:
\[
\mathcal P=\mathcal W^{-1}\circ\mathcal R\circ\mathcal M\circ\mathcal W\circ\mathcal C.
\]
\end{proposition}
Jeder Operator vertritt ein eigenes Gebiet: Koordinatensysteme, geometrische
Modellierung, physische Prozesse oder Datentransformation.

\subsection{Prinzip 5: Entwurf im realisierten Raum}
\begin{proposition}
Optimierungsziele müssen an realisierter Geometrie und nicht an theoretischen
Modellen bewertet werden.
\end{proposition}
Die maßgebliche Abbildung lautet
\[
\kappa_{\mathrm{target}}\mapsto\gamma_{\mathrm{real}}.
\]
Der Entwurf muss die durch Realisierung verursachte Transformation einbeziehen.

\subsection{Prinzip 6: Nutzung der Struktur}
\begin{proposition}
Die intrinsische Block- und Dünnbesetzungsstruktur des Problems muss für
effiziente Berechnung genutzt werden.
\end{proposition}
Dies umfasst Lokalität in der Bogenlänge, Blockzerlegung von Krümmung und
Überhöhung sowie dünn besetzte Jacobi-Matrizen.

\subsection{Prinzip 7: Maßstabsbewusstsein}
\begin{proposition}
Verschiedene Pipeline-Komponenten müssen auf geeigneten Maßstäben bewertet
werden.
\end{proposition}
CRS dominieren im großen, parametrische Modellierung im mittleren und
Korrekturen sowie Realisierung im kleinen Maßstab.

\subsection{Prinzip 8: Integration von Messungen}
\begin{proposition}
Messdaten müssen durch die World--Track-Transformation und nicht unmittelbar
im World-Raum integriert werden.
\end{proposition}
Dies gewährleistet Konsistenz mit der intrinsischen Alignment-Darstellung.

\subsection{Prinzip 9: Regularisierung als Notwendigkeit}
\begin{proposition}
Regularisierung ist ein wesentlicher Bestandteil von Alignment-Modellierung
und -Optimierung.
\end{proposition}
Sie stabilisiert inverse Probleme, hybride Korrekturen und numerische
Optimierung.

\subsection{Prinzip 10: Akzeptanz der Nicht-Eindeutigkeit}
\begin{proposition}
Alignment-Lösungen sind nicht eindeutig und müssen nach Engineering-Kriterien
bewertet werden.
\end{proposition}
Verschiedene Modelle können nahezu identische realisierte Geometrie erzeugen.
Lösungsqualität muss deshalb durch Dynamik, Bedingungen und Robustheit
bestimmt werden.

\subsection{Zusammenfassung}
\begin{summary}
Der AIM-Rahmen wird von Entwurfsprinzipien bestimmt, die aus dem
Zusammenwirken von Geometrie, Physik, Daten und Berechnung hervorgehen.

Sie definieren, wie Alignments dargestellt, transformiert und optimiert und
wie ihre Grenzen behandelt werden sollen. Gemeinsam bilden sie eine kohärente
Grundlage für Alignment-Based Information Modelling als einheitliche
Engineering- und rechnerische Disziplin.
\end{summary}
